% Further H9 deduction in run JAC1964-20261121.
% /root/h9_height; GPT-6 Astra Pro; effort setting not exposed.
% Classical results used below are explicitly cited; no formal verification
% or historical originality is asserted.

\subsection*{H9-T9. The algebraically trivial domain: incidence and trace}

Let $k$ be algebraically closed, let $C/k$ be a smooth projective connected
curve, and put $K=k(C)$. Let $Y/K$ be smooth projective of dimension
$2m-1$. Write $H_V$ for the degree of the unsigned geometric intersection
height on a smooth projective $K$-variety $V$; thus the signed middle
height on $Y$ is
\[
 h_m(x,y)=(-1)^m H_Y(x,y).
\]
We use the pairing of Kahn, \emph{Refined height pairing}, Theorem 4.14,
printed p.~31. Its restriction to algebraically trivial cycles exists in
every characteristic by Theorem 5.6(a), printed pp.~33--35. On a model
with admissible extensions it is their intersection degree, so this
normalization agrees with H9-T1.

Let $P=P^m(Y)$ be the rationally descended incidence Picard variety, and
let $a$ denote its quotient map on algebraically trivial cycles. Choose
an abelian direct summand $B\subset P$ in the isogeny category, and a
rational family $U\in\operatorname{CH}^m(B\times Y)_{\mathbf Q}$
whose family map is the inclusion $B\to P$. The existence over $K$ is
the rational Poincar\'e-family descent in JAC-P00. Subtract $U_0$ so
that the family is pointed. Let
\[
 D=\bigl(\operatorname{pr}_{12*}(U_{13}U_{23})\bigr)_{11},
 \qquad \ell=\frac{(-1)^m}{2}\Delta_B^*D
       \in\operatorname{NS}(B)_{\mathbf Q},                 \tag{H9.25}
\]
where the subscript removes the two axis restrictions. Suppose $\ell$
is a rational polarization. For the primitive summand of a generic
Lefschetz section of an even-dimensional abelian variety this hypothesis
is supplied by H9-T4, H9-T6, H9-T6a and JAC-P07, accepting 032.

Define
\[
 \mathcal A_B=
 \{z\in\operatorname{CH}^{m}(Y)_{\mathrm{alg},\mathbf Q}:
                      a(z)\in B(K)\otimes\mathbf Q\}.
\]
One may also use the group of $K$-cycles which become algebraically
trivial after a finite extension: admissibility and the curve-family
arguments below are obtained over that extension and descended by
transfer, with rational coefficients. The quotient map and the
subspace are understood in the rational category; they are independent
of clearing the denominators of the chosen isogenies.

\textbf{Claim.} The height on $\mathcal A_B$ is positive semidefinite,
and its radical is exactly
\[
 \{z\in\mathcal A_B:
      a(z)\in\operatorname{im}((\operatorname{Tr}_{K/k}B)(k)
                                  \longrightarrow B(K))
                                  \otimes\mathbf Q\}.       \tag{H9.26}
\]
More precisely, if $L$ is a symmetric ample line bundle representing
$N\ell$ for some positive integer $N$, and
\[
 \phi_L(b)=t_b^*L\otimes L^{-1}\in\operatorname{Pic}^0(B),
\]
then
\[
 h_m(z,z')=\frac1N\,
 H_B\bigl(\phi_L(a(z')),\,[a(z)]-[0]\bigr).                 \tag{H9.27}
\]
The right side uses the divisor--Albanese pairing, so the notation is
extended bilinearly to rational points. No assertion is made that an
arbitrary homologically trivial cycle has a point $a(z)$.
Explicitly, clear a positive denominator $r$ so that $rb$ is an actual
$K$-point; the notation $[b]-[0]$ in that pairing means the Albanese
class of $r^{-1}([rb]-[0])$. Its value is independent of $r$ by the
Albanese factorization. The same clearing of positive denominators
turns every rational inclusion, projector and family map used here
into an actual homomorphism before applying it to points.
The induced map from $\mathcal A_B$ modulo the subgroup in (H9.26) to
$\operatorname{LN}(B,K/k)\otimes\mathbf Q$ is an isomorphism with
bilinear forms: surjectivity is supplied by the pointed family $U$.

\paragraph{Proof: the required correspondence adjointness.}
Kahn's Proposition 3.8, printed pp.~19--20, and Theorem 4.14(b) state
adjointness for correspondences of degree zero. The degree shift needed
here reduces to exactly that statement. Given
$\Gamma\in\operatorname{CH}^{m}(T\times Y)_{\mathbf Q}$, a degree-zero
zero-cycle $\delta$ on $T$, and an algebraically trivial middle cycle
$\eta$ on $Y$, put $W=T\times\mathbf P^{m-1}$ and
\[
 \Gamma'=\Gamma\boxtimes H^{m-1}
   \in\operatorname{CH}^{2m-1}(W\times Y)_{\mathbf Q},
\]
where $H$ is the hyperplane class. This is a degree-zero
correspondence in Kahn's contravariant convention. Its transpose action
on $\eta$ is $\Gamma^t_*\eta\boxtimes H^{m-1}$, while its pushforward
action on $\delta\boxtimes[\mathbf P^{m-1}]$ is $\Gamma_*\delta$.
Adjointness, followed by adjointness for the projection $W\to T$, gives
\[
 H_Y(\Gamma_*\delta,\eta)
    =H_T(\delta,\Gamma^t_*\eta).                            \tag{H9.28}
\]
The projection formula uses only
$\deg_{\mathbf P^{m-1}}H^{m-1}=1$. Every cycle paired here is
algebraically trivial, so Kahn's admissibility theorem applies. This
argument also includes $m=1$, with $\mathbf P^0$.

\paragraph{Proof: the incidence kernel is killed by the height.}
Let $\eta$ have zero incidence image. An algebraically trivial cycle
$z$ is a finite sum of curve-family differences
$\Gamma_*\delta$, after a finite extension if necessary. By definition
of incidence equivalence, $\Gamma^t_*\eta$ is zero in
$\operatorname{Pic}^0(T)_{\mathbf Q}$ for every such tester. If this
vanishing is initially over $\overline K$, it already holds rationally
over the finite field of definition: the kernel of restriction on a
Picard group is killed by transfer after a suitable finite extension.
Equation (H9.28) therefore proves $H_Y(z,\eta)=0$. Finite base change
multiplies the degree height by its positive field degree, by Kahn's
Theorem 4.14(b) and Proposition 3.9(iv), so the assertion descends to
$K$.

Consequently the height on algebraically trivial cycles factors through
the incidence quotient. In particular, if $b=a(z)\in B(K)\otimes
\mathbf Q$, then $z$ may be replaced in the pairing by
$U_*([b]-[0])$. The difference has incidence image zero. The same
applies to $z'$.

Apply (H9.28) with $T=B$ and $\Gamma=U$. The composition $U^t\circ U$
is the divisor correspondence before rigidification in (H9.25).
Pullbacks of divisors from either factor act trivially on degree-zero
zero-cycles, so its action is that of $D$. The symmetric rigidified
divisor $D$ is the biextension divisor
\[
 D=\frac{(-1)^m}{N}
       (\mu^*L-p_1^*L-p_2^*L)
       \quad\text{in }\operatorname{Pic}(B\times B)_{\mathbf Q}
       \text{ after rigidification}.                        \tag{H9.29}
\]
Indeed symmetric divisor correspondences rigidified on the axes are
identified with symmetric homomorphisms $B\to\widehat B$, and
diagonal pullback is twice the associated N\'eron--Severi class.
Alternatively, the difference in (H9.29) has zero N\'eron--Severi class,
and every line bundle in $\operatorname{Pic}^0(B\times B)$ is a tensor
product of pullbacks from the factors; rigidification removes them.
The restriction of the displayed biextension at its first argument $b$
is precisely $\phi_L(b)$ under the stated translation convention.
This proves (H9.27), including its sign.

\paragraph{Proof: an explicit polarization--intersection identity.}
Let $d=\dim B>0$, and define the homomorphism
\[
 T_L:\operatorname{Pic}^0(B)\longrightarrow B,
 \qquad \alpha\longmapsto
        \operatorname{alb}_B(L^{d-1}\alpha).
\]
Intersection with a fixed cycle followed by the Albanese map defines
a homomorphism of abelian varieties. We claim the identity
\[
 T_L\phi_L=-c_L\,1_B,
 \qquad c_L=(d-1)!\chi(L)>0.                                 \tag{H9.30}
\]
We give an algebraic proof that fixes the sign in every characteristic.
First, the Albanese image of an intersection containing at least two
algebraically trivial divisor factors is zero. To see this, vary those
factors independently in copies of $\operatorname{Pic}^0(B)$, using
the Poincar\'e line bundles. Intersection followed by Albanese gives a
morphism from the product of these abelian varieties to $B$, taking
the origin to zero. Such a morphism is a homomorphism by rigidity, and
therefore is the sum of its restrictions to the factors. Each such
restriction is zero because another factor is the trivial line bundle.
The morphism is zero. This also proves the assertion when several
factors coincide, by restriction to a diagonal.

Now write $\alpha=\phi_L(b)$, so $t_b^*L=L+\alpha$ in additive
notation. Expanding the $d$th power and applying Albanese gives
\[
 \operatorname{alb}_B((t_b^*L)^d-L^d)=d\,T_L(\alpha).
\]
On the other hand, pullback by $t_b$ translates every point of the
zero-cycle $L^d$ by $-b$. Its Albanese sum therefore changes by
$-\deg(L^d)b=-d!\chi(L)b$, by Riemann--Roch for abelian varieties.
Dividing by $d$ proves (H9.30). Both sides are homomorphisms, so
checking on geometric points proves the identity over $K$ as well.

\paragraph{Proof: positivity and the exact radical.}
Kahn's Theorem 6.2, printed pp.~38--39, shows that the divisor--zero-cycle
height factors through the Albanese map and, after taking degree, through
the two Lang--N\'eron quotients. Its Theorem 6.6, printed pp.~41--42,
shows that
\[
 q_L(\alpha)=H_B(\alpha,L^{d-1}\alpha)
             =H_B(\alpha,T_L\alpha)
\]
is negative definite on
$\operatorname{LN}(\operatorname{Pic}^0(B),K/k)\otimes\mathbf Q$.
There is no conjectural hypothesis here, because the base $C$ has
dimension one. Since $\phi_L$ is an isogeny, it induces an isomorphism
of rational Lang--N\'eron groups. Equation (H9.30) gives
\[
 H_B(\phi_L(b),[b]-[0])
       =-\frac{1}{c_L}\,q_L(\phi_L(b)).                     \tag{H9.31}
\]
Thus this is positive definite modulo the trace, with exactly the trace
as its rational zero space. Equations (H9.27) and (H9.31) prove
semipositivity and (H9.26). If $B=0$, the incidence-kernel argument
already gives the same conclusion.\hfill$\square$

\paragraph{What this repairs, and what it does not assert.}
This is a precise trace-kernel theorem on the algebraically trivial
domain where the incidence quotient is defined. It proves the height
consequence of a primitive incidence polarization in every characteristic.
The proof identifies the pairing with a divisor--Albanese height and
derives its positivity directly; it does not invoke an unspecified
comparison theorem for a higher Abel--Jacobi map. It supplies no map
from every homologically trivial cycle to $P^m(Y)$, and therefore does
not repair the overlarge domain in literal H9(b). The E4 example in
JAC-P04 and the quadric example in H9-C1 retain their stated force.

\paragraph{Source and novelty record.}
The primary source used for the height properties is
\url{https://arxiv.org/pdf/2009.00533}, version 7, directly read in this
run: Proposition 3.8, pp.~19--20; Theorem 4.14, p.~31; Theorem 5.6(a),
pp.~33--35; Theorem 6.2, pp.~38--39; and Theorem 6.6, pp.~41--42.
The padding argument, the explicit proof of (H9.30), and their assembly
with the incidence quotient are deductions written in this run. No
historical originality, integral universal cycle, or formal proof check
is claimed.
